% Adenda v2 a la Introducción: impacto esperado, contenido exigido por nombre en
% las Instrucciones VIU y ausente en 01-introduction.tex (sellada).
\par\bigskip
\noindent\textbf{Impacto esperado.}
El impacto científico previsto es una regla de composición entre capas en la que
cada interfaz expone un predicado verificable con su dominio declarado, de modo
que la ausencia de garantía extremo a extremo deje de ser un supuesto implícito
y pase a ser una condición explícita que otros trabajos puedan intentar cerrar.
En lo industrial, el efecto directo está en la admisión de transportes
cooperativos: un criterio de \emph{wrench} realizable permite rechazar antes de
mover aquellas coaliciones que la cobertura de capacidad habría aceptado, y
convierte el fallo en una abstención declarada en lugar de un intento fallido en
planta. En lo social, la dimensión relevante es la seguridad de quien trabaja
junto a cargas pesadas movidas por varios vehículos: un sistema que se abstiene
cuando no puede certificar traslada la decisión a una persona en vez de
ejecutarla sin respaldo. Ninguno de los tres impactos se cuantifica en esta
memoria; se declaran como el resultado que el trabajo persigue y frente al cual
debe leerse su alcance.

\par\medskip
\noindent\textbf{Preguntas de investigación.}
El trabajo se articula en torno a seis preguntas, que la
Sección~\ref{subsec:conclusions-megajuego} y el apartado de respuestas del
Capítulo~\ref{sec:conclusions} cierran una a una.

\begin{enumerate}[label=\textbf{RQ\arabic*.},leftmargin=3.1em,itemsep=0.15em,parsep=0pt]
  \item ¿Bajo qué condiciones sobre recursos y penalización puede una regla
        local de formación de coaliciones regular cuotas exactas, y qué se
        pierde cuando esas condiciones no se cumplen? (OE1)
  \item ¿Acoplar el coste espacial a la decisión estratégica acelera la
        terminación, y sostiene la mensajería vecinal la ejecución cuando la
        red se degrada? (OE3)
  \item ¿Qué predicado certifica que una coalición puede ejercer físicamente el
        \emph{wrench} que la misión exige, y por qué no bastan la cardinalidad
        ni la cobertura de servicio? (OE2)
  \item ¿Qué fracción de los certificados perdidos por un fallo se recupera con
        reparación local, y en qué régimen deja de recuperarse? (OE3)
  \item ¿Cuánto cuesta en mensajes y cómputo coordinar rutas con reserva, y qué
        se gana frente a no reservar? (OE4, OE5)
  \item ¿Puede certificarse por bloque una arquitectura que encadene las tres
        capas, y sobrevive ese certificado a un motor de simulación distinto
        del que lo generó? (OE6)
\end{enumerate}
